% ID: SEL-D0032-OB-E04
% ESTADO: cotejada. Segunda pasada visual y borrador compilado revisados.
% Fuente: Selectividad/matematicas_aplicadas/titular_junio_EXAMEN 5-ANDALUCÍA.pdf
% SHA-256: 5f24b200fb6cf34f6d72016ec3dc35787b4074de9b7752677e5b0d959ed0dfcd
% Página: 2; opción B; ejercicio original 4.

La puntuación obtenida por los participantes en una prueba es una variable aleatoria que sigue una distribución Normal con una desviación típica de 6 puntos. Se toma una muestra aleatoria de 64 participantes en esa prueba, resultando una puntuación media de 35 puntos.
\begin{enumerate}[label=\alph*)]
\item \textbf{(1.25 puntos)} Calcule un intervalo de confianza, al 95\%, para la calificación media del total de participantes en la citada prueba.
\item \textbf{(1.25 puntos)} Halle el tamaño mínimo de la muestra necesaria para estimar la puntuación media del total de participantes, con un error inferior a 0.5 puntos y un nivel de confianza del 99\%.
\end{enumerate}
